\section{Problem Formulation}
\label{sec:model}

\subsection{Applications, Captures and States}

We view the application under test (AUT), as experienced by a fixed role $r$, as a
labelled transition system $M_r=(Q,A,\delta,q_0)$ with concrete UI configurations
$Q$, actions $A$ and a partial transition function $\delta:Q\times A\rightharpoonup
Q$. The auditor never sees $Q$ itself. It sees a \emph{capture}
$\kappa(q)=(u,d,t,s)$ made of the URL, the serialised DOM, the accessibility (AX)
tree and optionally a screenshot. One role is fixed per run, so auditing several
roles means several runs.

\begin{definition}[State fingerprint]
Let $N$ normalise a capture by removing volatile content (session and CSRF
tokens, nonces, long numeric timestamps, generated hexadecimal identifiers,
per-session CDP node identifiers, transient focus) and let $H$ be SHA-256. The
fingerprint is $\phi=H\circ N$, and two configurations are the same abstract state
iff $\phi(\kappa(q))=\phi(\kappa(q'))$.
\end{definition}

A normaliser can be too sensitive, so that every page load is a new node, or too
coarse, so that a security-relevant change merges into an existing state and is
never audited. Sect.~\ref{sec:fp} probes both with targeted edits; it is a probe,
not a measurement of either failure rate.

\subsection{The Exploration Graph Is Cyclic}

\begin{definition}[Exploration graph]
The exploration graph is a directed multigraph $G=(V,E)$ with $V\subseteq
\operatorname{img}\phi$ and $E\subseteq V\times A\times V$. An edge $(v,a,v')$
records that executing $a$ in $v$ was observed to lead to $v'$. We call it a
\emph{revisit edge} if $v'$ had already been discovered when the edge was added.
\end{definition}

A revisit edge is any non-tree edge of the discovery order, including cross and
forward edges; it is not the depth-first notion of a back edge, although the
prototype's field is named \texttt{is\_back\_edge}. Crawlers often drop revisits as
redundant. We keep every edge, because facts such as ``log out leads to login''
describe the AUT's navigation structure and a path-sensitive policy can only be
checked on a graph that retains them. The graph is therefore cyclic and revisit
edges are never pruned. Loop prevention rests instead on a \emph{claim set}
$C\subseteq V\times A$ of visited (state, action) pairs.

\begin{definition}[Claim-before-execute]
A frontier action $a$ is executed from $v$ only if $(v,a)\notin C$ when it is
enqueued. The pair is inserted into $C$ before $a$ is sent to the browser, and the
edge $(v,a,v')$ is persisted only after $a$ has succeeded.
\end{definition}

One browser can only act on the page it is on, so a frontier action from
another state is reached by \emph{replaying} its path from the seed, which fires
the path's actions again. \sys{} limits this by action kind: only
$\textsf{ReplaySafe}=\{\texttt{navigate},\texttt{click},\texttt{back}\}$ may be
re-fired. A candidate whose path contains a \texttt{fill}, \texttt{select} or
\texttt{submit} is recorded as skipped with a reason, and its subtree is reported
as unexplored. The residual assumption is that clicks are navigational: a click
that mutates state on a replayed path, such as a one-click ``Delete'', \emph{is}
re-fired, which we record as an open limitation. Actions are content-addressed,
$\mathit{id}(a)=H(\mathit{role}\,\|\,\mathit{name}\,\|\,\mathit{selector})$, so one
``Log out'' control has one identity on every page, while deduplication is scoped
to the pair $(v,a)$.

\begin{proposition}[Bounded frontier executions, uninterrupted runs]
\label{prop:term}
Assume $|V|\le n$ (enforced by a state cap), at most $k$ actions per state, a
depth bound $D$ on explored states, and a run without crashes. Then the loop
terminates; it executes at most $nk$ frontier actions; no pair $(v,a)$ is executed
as a frontier action twice; and replay re-fires at most $nkD$ actions in total.
\end{proposition}
\begin{proof}
Each frontier execution first adds a new pair to $C$, and $|C|\le nk$. A pair is
enqueued only when its state is first visited and the pair is not yet claimed or
queued, so it enters the frontier at most once. Every loop iteration removes one
frontier entry, so the loop halts. A replay re-fires the path of one frontier
entry, whose length equals the depth of its source state ($\le D$), and there are
at most $nk$ entries.
\end{proof}

The proposition is deliberately narrower than ``at-most-once'': it bounds
\emph{frontier} executions of a pair, not the times an action reaches the AUT, since
replay re-fires path actions (Sect.~\ref{sec:rq1} measures this), and it excludes
crashes. The intended behaviour is that a pair claimed but lacking an edge counts
as done and is not fired again, at the price of leaving it unexplored;
Sect.~\ref{sec:rq1} reports a simulated test in which two resume mechanisms
disagree about this.

\subsection{Policies as Two Set Operations}

A QA engineer writes the policy in English, to be compiled into clauses
$c=(\sigma_c,\rho_c,\pi_c)$: a subject $\sigma_c$, a role scope $\rho_c$ (empty
meaning all roles) and a polarity $\pi_c\in\{\textsf{must\_exist},\textsf{must\_not\_exist}\}$.
The clauses split into a forbidden set $\Forb$ and a required set $\Req$.
Perception maps each state $v$ to a set $\Sig(v)$ of grounded subjects: element
tags, accessible names and capabilities. For role $r$, $\Forb_r$ and $\Req_r$
denote the clauses whose scope admits $r$.

\begin{definition}[Negation-aware audit]
\begin{align}
\mathit{Viol}_\Forb(v) &= \{\,c\in\Forb_r : \sigma_c\in\Sig(v)\,\}
  && \text{(forbidden present)}\label{eq:forb}\\
\mathit{Viol}_\Req(v)  &= \{\,c\in\Req_r : \sigma_c\notin\Sig(v)\,\}
  && \text{(required absent)}\label{eq:req}
\end{align}
The audit of $v$ is $\mathit{Viol}_\Forb(v)\cup\mathit{Viol}_\Req(v)$, with each
violation tagged by the operation that produced it.
\end{definition}

Equation~(\ref{eq:forb}) is an intersection and Eq.~(\ref{eq:req}) a set
difference. An oracle built from intersection alone cannot flag a missing element,
since a missing element has nothing to intersect with.

\subsection{Error Direction Under Imperfect Perception}
\label{sec:asym}

In practice $\Sig(v)$ is replaced by an estimate $\hat\Sig(v)$ produced by a
DOM/AX correlator or a VLM. The two operations respond to its errors in opposite
directions, and this follows from their definitions.

\begin{proposition}[Monotonicity of the two operations]
\label{prop:asym}
Let $\widehat{\mathit{Viol}}$ be the audit computed from $\hat\Sig$.
\begin{enumerate}
\item If perception under-reports, $\hat\Sig(v)\subseteq\Sig(v)$, then
$\widehat{\mathit{Viol}}_\Forb(v)\subseteq\mathit{Viol}_\Forb(v)$ and
$\widehat{\mathit{Viol}}_\Req(v)\supseteq\mathit{Viol}_\Req(v)$: forbidden clauses
can only be missed and required clauses can only raise false alarms.
\item If perception over-reports, $\hat\Sig(v)\supseteq\Sig(v)$, both inclusions
are reversed.
\end{enumerate}
\end{proposition}
\begin{proof}
Equation~(\ref{eq:forb}) is monotone in $\Sig$ and Eq.~(\ref{eq:req}) is antitone
in $\Sig$.
\end{proof}

This is an immediate property of the definitions, not a research result. Its use
is a testable prediction about an implementation: for a perceiver that only drops
elements, the engine must produce no forbidden false positives and no required
false negatives, and for one that only adds spurious elements, no forbidden false
negatives and no required false positives. A violation would mean the engine
deviates from Eqs.~(\ref{eq:forb}) and~(\ref{eq:req}), for example through
substring matching; Sect.~\ref{sec:rq2} tests this on the production engine. A
perceiver with a single global bias cannot be safe for both clause types, although
one with polarity-specific bias could be.

For real channels the direction is relative to the policy vocabulary: a channel
that tags an element with a string outside the policy's subject vocabulary has,
for that clause, under-reported. That reading makes the direction checkable on a
rule-based channel but also means one application cannot refute the observation;
the perturbation test is the falsifiable part. Which polarity a VLM errs towards on
UI screenshots
\cite{alhamoud2025negation,yuksekgonul2023aro,ma2023crepe} is an empirical
question we leave to the probe set of Sect.~\ref{sec:discussion}.


